% arXiv submission: RP-06
% Title: The Fourfold Correspondence: Grammar, Automaton, Type, and Proof
%        from Chomsky to Curry-Howard-Lambek
% Author: A. Padavano (Independent researcher)
% Primary category: cs.FL (Formal Languages and Automata Theory)
% Cross-list: cs.LO, cs.CL, math.CT
% MSC 2020: 68Q45, 03B40, 03F03, 18C50
% License: CC BY 4.0
% Date: 2026

\documentclass[12pt,a4paper]{article}

\usepackage[utf8]{inputenc}
\usepackage[T1]{fontenc}
\usepackage{amsmath,amssymb,amsthm}
\usepackage{hyperref}
\usepackage[sort&compress,numbers]{natbib}
\usepackage{booktabs}
\usepackage{graphicx}
\usepackage{alltt}
\usepackage{microtype}
\usepackage[margin=1in]{geometry}

\hypersetup{
  colorlinks=true,
  linkcolor=blue,
  citecolor=blue,
  urlcolor=blue
}

\newtheorem{definition}{Definition}
\newtheorem{theorem}{Theorem}
\newtheorem{conjecture}{Conjecture}
\newtheorem{proposition}{Proposition}

\title{The Fourfold Correspondence:\\Grammar, Automaton, Type, and Proof\\from Chomsky to Curry-Howard-Lambek}

\author{A.~Padavano\\
  \textit{Independent researcher, Studium Generale ORGANVM}}

\date{2026}

\begin{document}

\maketitle

\begin{abstract}
This paper establishes a fourfold structural correspondence linking formal
grammars, abstract automata, type systems, and proof calculi. The classical
result in formal language theory pairs each level of the Chomsky hierarchy with
a class of abstract machines. We extend this twofold correspondence to a
fourfold one by demonstrating that each level admits a canonical
type-theoretic characterisation and a corresponding proof-theoretic regime.
The extension is mediated by categorial grammar (Ajdukiewicz, Bar-Hillel,
Lambek) and the Curry-Howard-Lambek correspondence. We argue that these are
not four independent correspondences but four projections of a single
mathematical structure, while carefully stratifying evidence strength at each
level: the correspondence ranges from constructive isomorphism at the
context-free level, through structural embedding at the context-sensitive
level, to analogy at the unrestricted level. We ground linguistic claims in
typologically diverse evidence spanning agglutinative (Turkish),
isolating/tonal (Mandarin), sign (ASL), non-configurational (Warlpiri), and
polysynthetic (Mohawk) languages, addressing limitations posed by gradient
grammaticality, multiword expressions, and contested empirical arguments. We
survey the historical development from Chomsky's 1956 hierarchy through
Lambek's 1958 syntactic calculus to the DisCoCat framework of Coecke,
Sadrzadeh, and Clark, examining where the fourfold correspondence holds
precisely and where it breaks down. We identify three open problems: the
type-theoretic characterisation of mildly context-sensitive languages (for
which we propose a conjecture), whether neural language models learn
type-theoretic structure, and the formal relationship between Chomsky's Merge
operation and typed function application.
\end{abstract}

\noindent\textbf{Keywords:} Chomsky hierarchy, type theory, Curry-Howard
correspondence, categorial grammar, Lambek calculus, formal language theory,
mildly context-sensitive languages, compositional semantics

\bigskip

\section{Introduction}

In 1956, Noam Chomsky proposed a classification of formal grammars into four
types, arranged by increasing generative power, that has since become one of
the most consequential taxonomies in the mathematical sciences~\cite{Chomsky1956}.
The \emph{Chomsky hierarchy}, as it came to be known, stratifies formal
grammars into regular grammars (Type~3), context-free grammars (Type~2),
context-sensitive grammars (Type~1), and unrestricted grammars (Type~0). Each
class generates a strictly larger family of formal languages than its
predecessor, and each admits recognition by a correspondingly more powerful
class of abstract machines: finite automata, pushdown automata, linear bounded
automata, and Turing machines, respectively. This grammar-automaton
correspondence, established through the combined efforts of
Chomsky~\cite{Chomsky1959}, Rabin and Scott~\cite{RabinScott1959},
Myhill~\cite{Myhill1957}, and others, constitutes one of the foundational
results of theoretical computer science.

The significance of the hierarchy extends well beyond its original domain.
Chomsky formulated it to characterise the structural properties of natural
language; within a decade, the same formal apparatus had been adopted by
computer scientists as the basis for programming language definition. Backus
and Naur's notation for the syntax of ALGOL~60~\cite{Backus1959,Naur1960} is
a context-free grammar in Chomsky's precise sense, and the phrase-structure
rules that Chomsky proposed for English in \emph{Syntactic
Structures}~\cite{Chomsky1957} are formally identical to the production rules
that define programming language parsers. The hierarchy thus serves as a shared
substrate linking linguistics and computer science---a mathematical Rosetta
Stone.

Yet the grammar-automaton correspondence, for all its elegance, captures only
two of the relevant dimensions. There exists a third dimension---type-theoretic---that is less widely recognised but arguably runs at least as deep. In
this dimension, each level of the Chomsky hierarchy corresponds to a regime of
type-system expressiveness: regular grammars to finite (enumeration) types,
context-free grammars to recursive algebraic types, context-sensitive grammars
to parametrically polymorphic types, and unrestricted grammars to fully
dependent types. This correspondence can be made precise through the framework
of \emph{categorial grammar}, in which syntactic categories are function
types, grammatical derivations are proofs, and sentence formation is typed
function application~\cite{Ajdukiewicz1935,BarHillel1953,Lambek1958}.

The bridge from types to proofs is provided by the \emph{Curry-Howard
correspondence}~\cite{Howard1980}, which demonstrates that propositions are
types and proofs are programs. Lambek~\cite{Lambek1972,Lambek1980} extended
this to a three-way equivalence by showing that typed lambda calculi
correspond to Cartesian closed categories. When combined with the observation
that Lambek's own 1958 syntactic calculus is simultaneously a grammar
formalism, a logical system, and a typed calculus, the three-way
correspondence becomes fourfold: grammatical derivations are proofs are typed
terms are morphisms in a category.

\subsection{Stratification of Correspondence Strength}

A critical methodological point must be established at the outset. The nature
of the fourfold correspondence varies significantly across the levels of the
Chomsky hierarchy, and conflating these different degrees of strength would
undermine the contribution's formal credibility. We distinguish three tiers:

\textbf{Isomorphism.} A constructive bijection exists between the structures
on both sides, with inverse maps that preserve all relevant structure.

\textbf{Structural embedding (homomorphism).} A structure-preserving map
exists from one side to the other, but the map is not necessarily surjective or
invertible.

\textbf{Analogy.} The two sides share structural features---similar closure
properties, similar decidability boundaries, similar compositional
principles---but no single formal map has been established.

\begin{table}[ht]
\centering
\caption{Stratification of correspondence strength across hierarchy levels}
\label{tab:stratification}
\begin{tabular}{@{}lllll@{}}
\toprule
Chomsky Level & Grammar-Automaton & Grammar-Type & Grammar-Proof & Overall \\
\midrule
Type 3 (Regular) & Isomorphism & Embedding & Embedding & Embedding \\
Type 2 (Context-free) & Isomorphism & \textbf{Isomorphism} & Isomorphism & \textbf{Isomorphism} \\
MCS & Isomorphism & \textbf{Open} & Partial & Partial \\
Type 1 (Context-sensitive) & Isomorphism & Embedding (conj.) & Analogy & Embedding \\
Type 0 (Unrestricted) & Isomorphism & Analogy & Analogy & Analogy \\
\bottomrule
\end{tabular}
\end{table}

This stratification has a crucial consequence: the paper's strongest
contribution is at the context-free level, where the fourfold correspondence
is a genuine isomorphism. Claims about the other levels must be qualified
accordingly, and we do so throughout.

\subsection{Scope and Framing}

This paper is framed within the generativist tradition of formal grammar, for
the reason that the Chomsky hierarchy \emph{is} a generativist construct. We
acknowledge at the outset that substantial traditions in linguistics---including
Construction Grammar~\cite{Goldberg1995,Goldberg2006}, Cognitive
Grammar~\cite{Langacker1987,Langacker2008}, and usage-based
approaches~\cite{Bybee2006,Bybee2010}---reject the premise that natural
language structure is best described by formal grammars in the Chomsky sense.
Section~\ref{sec:limitations} addresses this debate directly. The fourfold
correspondence applies most cleanly to \emph{formal} and \emph{computational}
languages, where compositionality is exact and grammaticality is binary. Its
extension to natural language is conjectural and approximate, and we are
explicit about where the fit degrades.

Three research questions guide the inquiry:

\textbf{RQ1.} What is the precise type-theoretic characterisation of each
level of the Chomsky hierarchy, and can this characterisation be extended to
the class of \emph{mildly context-sensitive} languages---the conjectured
complexity class of natural language~\cite{Joshi1985}?

\textbf{RQ2.} Do neural language models (transformer architectures) implicitly
learn type-theoretic structure, and if so, at what level of the Chomsky
hierarchy does their implicit grammar reside?

\textbf{RQ3.} Is the Merge operation of Chomsky's Minimalist
Program~\cite{Chomsky1995} formally equivalent to function application in the
typed lambda calculus?

\section{The Chomsky Hierarchy}

The hierarchy classifies formal grammars into four nested types, each defined
by progressively weaker constraints on the form of production rules. A formal
grammar is a quadruple $G = (V, \Sigma, R, S)$, where $V$ is a finite set of
nonterminal symbols, $\Sigma$ is a finite set of terminal symbols disjoint
from $V$, $R$ is a finite set of production rules, and $S \in V$ is a
distinguished start symbol.

\subsection{Type 3: Regular Languages}

\textbf{Grammar restriction.} A regular grammar restricts production rules to
the forms $A \to aB$ or $A \to a$, where $A, B \in V$ and $a \in \Sigma$.

\textbf{Formal languages generated.} Regular grammars generate exactly the
class of \emph{regular languages}. Kleene's theorem~\cite{Kleene1956}
established the equivalence of regular expressions and finite automata. Regular
languages are closed under union, intersection, complementation, concatenation,
and Kleene star.

\textbf{Linguistic examples.} Turkish vowel harmony provides a well-studied
example from a typologically distinct, agglutinative language. Vowel
harmony---the requirement that vowels in suffixes agree in backness and
rounding with the stem vowel---can be modelled by a finite-state
transducer~\cite{KaplanKay1994}. Mandarin Chinese tone sandhi offers an
example from an isolating, tonal language: third-tone sandhi is a local,
context-dependent phonological process capturable by finite-state rules.

\textbf{Type-theoretic analogue.} The type-theoretic counterpart of a regular
language is a \emph{finite type} or \emph{enumeration type}. The states of a
finite automaton correspond to the variants of an enumeration type, and the
transition function corresponds to a case-analysis (pattern match). The
correspondence at this level is a \emph{structural embedding} rather than an
isomorphism.

\subsection{Type 2: Context-Free Languages}

\textbf{Grammar restriction.} A context-free grammar restricts production
rules to the form $A \to \alpha$, where $A$ is a single nonterminal and
$\alpha$ is any string of terminals and nonterminals.

\textbf{The recognition machine: pushdown automata.} A language is
context-free if and only if it is accepted by some nondeterministic pushdown
automaton~\cite{Chomsky1962,Evey1963}.

\textbf{Linguistic examples.} Warlpiri, a non-configurational Australian
language~\cite{Hale1983}, challenges the context-free model: its free word
order with discontinuous constituents cannot be directly generated by standard
CFGs. This demonstrates that while CFGs capture the \emph{hierarchical
dependency} structure, they do not capture \emph{linear realisation} in free
word-order languages.

\textbf{Type-theoretic analogue: recursive algebraic types.} This
correspondence is exact. A context-free grammar is a system of equations
defining nonterminals in terms of terminals and other nonterminals, possibly
recursively. A recursive algebraic data type is a system of equations defining
types in terms of base types and other type constructors, possibly recursively.
The BNF rule:

\begin{verbatim}
Expr ::= Expr '+' Term | Term
Term ::= Term '*' Factor | Factor
Factor ::= '(' Expr ')' | Number
\end{verbatim}

\noindent is isomorphic to the Haskell type definition:

\begin{verbatim}
data Expr   = Add Expr Term | ETerm Term
data Term   = Mul Term Factor | TFactor Factor
data Factor = Parens Expr | Num Int
\end{verbatim}

The nonterminals of the grammar correspond to the type names; alternatives
separated by \texttt{|} correspond to sum types; sequencing corresponds to
product types. Parsing a string is exactly constructing a value of the
corresponding algebraic type. \emph{Strength: isomorphism.}

\subsection{Type 1: Context-Sensitive Languages}

\textbf{Grammar restriction.} A context-sensitive grammar permits production
rules of the form $\alpha A \beta \to \alpha \gamma \beta$, where $A$ is a
nonterminal and $\gamma$ is non-empty.

\textbf{Linguistic examples.} Swiss German cross-serial
dependencies~\cite{Shieber1985} generate patterns of the form
$\{a^n b^n c^n\}$. We note that Shieber's argument, while widely accepted, is
not uncontested~\cite{PullumGazdar1982,MichaelisKracht1997}. Mohawk noun
incorporation~\cite{Baker1988,Baker1996} presents a polysynthetic challenge.

\textbf{A worked example: encoding $\{a^n b^n c^n\}$ in System F.} The
string $a^i b^j c^k$ can be represented by a triple of Church numerals
$(\underline{i}, \underline{j}, \underline{k})$. The predicate $i = j = k$
can be verified in System~F using the standard encoding of Church numeral
equality. However, this encoding operates at the level of \emph{refinement
types} rather than pure System~F, and is not \emph{structural} in the way
that the BNF-to-algebraic-type correspondence is. We therefore classify the
Type~1 correspondence as a \emph{structural embedding}: context-sensitive
phenomena can be faithfully represented but the representation is an encoding
rather than an isomorphism.

\subsection{Type 0: Recursively Enumerable Languages}

Unrestricted grammars generate exactly the \emph{recursively enumerable}
languages. The type-theoretic counterpart is \emph{full dependent type
theory}: the Calculus of Constructions~\cite{CoquandHuet1988} and
Martin-L\"{o}f's Intuitionistic Type Theory~\cite{MartinLof1975,MartinLof1984}.
In their full generality, dependent type systems are Turing-complete:
type-checking is equivalent to theorem-proving, which is undecidable.
\emph{Strength: analogy.}

\subsection{Beyond the Hierarchy: Mildly Context-Sensitive Languages}

Between the context-free and context-sensitive levels lies a family of language
classes termed \emph{mildly context-sensitive}, introduced by
Joshi~\cite{Joshi1985}. Several independently developed grammar formalisms
satisfy Joshi's criteria: tree-adjoining grammars (TAGs;
\cite{JoshiLevyTakahashi1975}), combinatory categorial grammars
(CCGs; \cite{Steedman2000}), head grammars~\cite{Pollard1984}, and linear
indexed grammars~\cite{Gazdar1988}. These were shown to be weakly
equivalent~\cite{VijayShankerWeir1994}.

American Sign Language (ASL) challenges the string-based foundation of the
hierarchy: ASL uses \emph{simultaneous} articulation across multiple channels
(two hands, facial expression, eye gaze, body
posture)~\cite{SandlerLilloMartin2006}.

\subsection{Summary Table}

\begin{table}[ht]
\centering
\caption{Grammar-automaton-type correspondence across all levels}
\label{tab:hierarchy}
\small
\begin{tabular}{@{}llllll@{}}
\toprule
Type & Grammar & Automaton & Type Analogue & Decidability & Strength \\
\midrule
3 & Regular & Finite automaton & Finite types & $O(n)$ & Embedding \\
2 & Context-free & Pushdown automaton & Recursive algebraic & $O(n^3)$ & \textbf{Isomorphism} \\
MCS & TAG/CCG & Embedded PDA & Bounded 2nd-order (conj.) & $O(n^6)$ & Partial \\
1 & Context-sensitive & Linear bounded & Parametric polymorphism & PSPACE & Embedding \\
0 & Unrestricted & Turing machine & Full dependent types & Undecidable & Analogy \\
\bottomrule
\end{tabular}
\end{table}

\section{From Grammars to Machines}

The grammar-automaton correspondence is established by constructive
translations between grammar types and automaton types at each
level~\cite{Kleene1956,Chomsky1962,Kuroda1964}. The computational complexity
of the membership problem varies dramatically:

\begin{itemize}
  \item \textbf{Regular:} $O(n)$ for deterministic finite automata.
  \item \textbf{Context-free:} $O(n^3)$ for CYK; $O(n)$ for LR($k$) grammars.
  \item \textbf{Context-sensitive:} PSPACE-complete.
  \item \textbf{Recursively enumerable:} Undecidable.
\end{itemize}

The sharp jump from polynomial to PSPACE-complete explains why programming
languages are designed to have context-free syntax.

Natural languages cluster at the mildly context-sensitive level, for reasons
potentially involving learnability~\cite{Joshi1985}, processing
constraints~\cite{Stabler2013}, and the syntax-semantics
interface~\cite{Steedman2000}.

\section{From Grammars to Types}

\subsection{The Lambek Calculus: Grammar as Logic}

Lambek's \emph{syntactic calculus}~\cite{Lambek1958} reformulates categorial
grammar as a logical system. It operates over basic types ($S$, $N$, $NP$) and
two type-forming connectives: right division $A/B$ and left division
$B \backslash A$. The calculus includes introduction and elimination rules for
both connectives, precisely analogous to natural deduction rules for
implication. The Lambek calculus is a \emph{non-commutative substructural
logic}, and Pentus~\cite{Pentus1997} proved that its product-free fragment
generates exactly the context-free languages.

\subsection{Categorial Grammar: Ajdukiewicz, Bar-Hillel, Lambek}

In categorial grammar, every word is assigned syntactic types, and sentences
are derived by functional application alone. The critical observation is that
categorial types \emph{are} function types in a typed lambda calculus.
The category $NP/N$ is the function type $N \to NP$. Sentence formation
\emph{is} function application. Parsing \emph{is} type-checking.

\subsection{Combinatory Categorial Grammar (Steedman)}

CCG~\cite{Steedman2000} extends classical categorial grammar with composition
($B$), type-raising ($T$), and other combinators corresponding to combinatory
logic~\cite{CurryFeys1958}. CCG generates exactly the mildly context-sensitive
languages~\cite{VijayShankerWeir1994}, making it arguably the most successful
realisation of the grammar-type-proof correspondence in computational practice.

\subsection{Montague Grammar and the Lambda Calculus Connection}

Montague's programme~\cite{Montague1970a,Montague1970b,Montague1973} forged
the definitive link between natural language grammar and the typed lambda
calculus. In category-theoretic terms, Montague grammar defines a
\emph{functor} from the syntactic category to the semantic category.

\subsection{The Grammar-Type Hierarchy}

\begin{itemize}
  \item \textbf{Regular $\leftrightarrow$ finite types.} Strength: embedding.
  \item \textbf{Context-free $\leftrightarrow$ recursive algebraic types.}
    Strength: \textbf{isomorphism}.
  \item \textbf{Context-sensitive $\leftrightarrow$ parametric polymorphism
    (System~F).} Strength: embedding (conjectured).
  \item \textbf{Unrestricted $\leftrightarrow$ full dependent types.}
    Strength: analogy.
\end{itemize}

\begin{conjecture}[MCS Gap Conjecture]
The type system corresponding to mildly context-sensitive languages is a
\emph{bounded linear polymorphic calculus}: a fragment of the linear lambda
calculus with second-order quantification restricted to bounded polymorphism
(System $F_{<:}$ restricted to linear types) or to a finite rank.
\end{conjecture}

This conjecture is falsifiable: it predicts that the languages typable in such
a system coincide with the MCS class.

\subsection{Pregroup Grammars}

Lambek's pregroup grammars~\cite{Lambek1999,Lambek2008} are a radical
algebraic simplification, where syntactic types are elements of a pregroup
with left and right adjoints satisfying $a^l \cdot a \leq 1 \leq a \cdot a^l$
and $a \cdot a^r \leq 1 \leq a^r \cdot a$.

\section{The Curry-Howard-Lambek Triangle and the Fourfold Correspondence}

\subsection{Proofs = Programs (Curry-Howard)}

The Curry-Howard correspondence establishes an isomorphism between formal
logic and typed lambda calculi:

\begin{table}[ht]
\centering
\caption{The Curry-Howard dictionary}
\label{tab:curryhoward}
\begin{tabular}{@{}ll@{}}
\toprule
Logic & Computation \\
\midrule
Proposition & Type \\
Proof & Program (term) \\
Implication $A \Rightarrow B$ & Function type $A \to B$ \\
Conjunction $A \wedge B$ & Product type $A \times B$ \\
Disjunction $A \vee B$ & Sum type $A + B$ \\
Universal $\forall x.\, P(x)$ & Dependent function $\Pi(x{:}A).\, P(x)$ \\
Existential $\exists x.\, P(x)$ & Dependent pair $\Sigma(x{:}A).\, P(x)$ \\
Modus ponens & Function application \\
Proof normalisation & $\beta$-reduction \\
\bottomrule
\end{tabular}
\end{table}

\subsection{Programs = Morphisms (Lambek)}

Lambek~\cite{Lambek1972,Lambek1980} demonstrated that typed lambda calculi
correspond to Cartesian closed categories, yielding:
\[
\text{Intuitionistic propositional logic}
  \cong \text{Simply typed }\lambda\text{-calculus}
  \cong \text{Cartesian closed categories}
\]

\subsection{The Linguistic Extension}

In categorial grammar, \emph{parsing a sentence is constructing a proof is
type-checking a term is composing morphisms in a category}. These are the same
activity described in four mathematical languages.

\begin{table}[ht]
\centering
\caption{The fourfold correspondence}
\label{tab:fourfold}
\begin{tabular}{@{}llll@{}}
\toprule
Grammar & Logic & Computation & Category \\
\midrule
Syntactic category & Proposition & Type & Object \\
Derivation & Proof & Program & Morphism \\
$A/B$ & Implication & Function type & Exponential \\
Application & Modus ponens & $\beta$-reduction & Evaluation \\
Lexicon & Axioms & Constants & Generators \\
Grammaticality & Provability & Type inhabitedness & Morphism existence \\
\bottomrule
\end{tabular}
\end{table}

\subsection{The Consolidated Fourfold Table}

\begin{table}[ht]
\centering
\caption{Full four-column, five-row correspondence}
\label{tab:consolidated}
\small
\begin{tabular}{@{}llllll@{}}
\toprule
Level & Grammar & Automaton & Type System & Proof Calculus & Strength \\
\midrule
Type 3 & Regular & Finite aut. & Finite types & Propositional & Embedding \\
Type 2 & CFG / Lambek & PDA & Recursive alg. types & Intuit. prop. logic & \textbf{Iso.} \\
MCS & TAG / CCG & Embedded PDA & Bounded lin. poly. (conj.) & Restricted 2nd-order & Partial \\
Type 1 & CSG & LBA & System F / Refinement & 2nd-order intuit. & Embedding \\
Type 0 & Unrestricted & TM & Full dep. types (CoC) & Higher-order logic & Analogy \\
\bottomrule
\end{tabular}
\end{table}

\subsection{Where It Breaks Down}
\label{sec:limitations}

\textbf{Gradient grammaticality.} The Chomsky hierarchy assumes binary
grammaticality, but natural language judgments are
gradient~\cite{LauClarkLappin2017,SprouseSchutzeAlmeida2013}. Graded monads
and quantitative type theories~\cite{Atkey2018,McBride2016} provide possible
extensions.

\textbf{The multiword expression problem.} A substantial proportion of natural
language (30--80\% of running text~\cite{ErmanWarren2000,Wray2002}) consists
of prefabricated chunks that resist compositional analysis. The fourfold
correspondence applies to the \emph{compositional} portion and must be
supplemented for the non-compositional portion.

\textbf{The generativist vs.\ usage-based debate.} Usage-based
approaches~\cite{Goldberg1995,Langacker1987,Bybee2006} reject the premise
that abstract formal grammars are the right description of linguistic
knowledge. A moderate usage-based position can potentially be accommodated by
extending the correspondence to weighted grammars and graded type systems.

\subsection{The DisCoCat Framework}

The DisCoCat framework~\cite{CoeckeSadrzadehClark2010} resolves the tension
between distributional and compositional semantics using category theory. A
DisCoCat model is a strong monoidal functor $F \colon \mathbf{Preg} \to
\mathbf{FdVect}$. The connection to quantum computing has led to
implementations on quantum hardware via the \texttt{lambeq}
framework~\cite{Kartsaklis2021}.

\section{Modern Developments}

\subsection{Programming Language Theory}

Bidirectional type-checking~\cite{DunfieldKrishnaswami2021} mirrors the
bidirectionality of parsing. Algebraic effects extend type classification from
values to computations. Homotopy Type Theory (HoTT) adds a geometric
interpretation: types are spaces, proofs are paths.

\subsection{NLP: CCG Parsing and Neural Approaches}

CCG parsers~\cite{ClarkCurran2007,LewisSteedman2014} achieve broad-coverage
parsing producing typed semantic representations. Probing studies of
transformers~\cite{HewittManning2019,Clark2019,Jawahar2019} provide
suggestive evidence that neural models encode at least context-free and some
context-sensitive structure, though the probing methodology has been
challenged~\cite{HewittLiang2019}.

\subsection{The Minimalist Program and Merge}

Chomsky's Merge operation~\cite{Chomsky1995,Chomsky2000} takes two syntactic
objects and combines them into a set. The parallel with function application
has been noted~\cite{Collins2002,Stabler2011}, and Minimalist Grammars
generate exactly the mildly context-sensitive
languages~\cite{Stabler1997,Stabler2011}. The key obstacle to full
equivalence is the symmetry of Merge vs.\ the asymmetry of application.

\section{Implications}

\subsection{For Programming Language Design}

The grammar hierarchy provides a decision framework for type system design:
finite types (trivial checking), recursive algebraic types (decidable),
parametric polymorphism (decidable checking, undecidable inference), dependent
types (decidable only with termination enforcement).

\subsection{For NLP}

The type-theoretic perspective suggests that parsing should be understood as
\emph{type inference}: producing typed lambda terms rather than untyped parse
trees. Type-driven parsing~\cite{BaldridgeKruijff2003} reduces ambiguity by
filtering type-inconsistent derivations.

\subsection{For Governance System Design}

The Chomsky hierarchy classifies the decidability of governance rule systems:
regular rules ($O(n)$, guaranteed termination), context-free rules
(polynomial, decidable), context-sensitive rules (PSPACE-complete),
unrestricted rules (potentially non-terminating). Governance designers should
choose the weakest level that suffices for their requirements.

\section{Discussion}

The fourfold correspondence admits three interpretations of varying strength:
the \emph{weak} (structural analogies), the \emph{moderate} (precise
mathematical maps at certain levels, conjectural at others), and the
\emph{strong} (four descriptions of a single mathematical structure). We have
presented evidence for the moderate interpretation while arguing that the
strong interpretation is a productive research programme. The primary gap is
at the mildly context-sensitive level.

The limitations of gradient grammaticality, multiword expressions, and the
generativist-vs.-usage-based debate are not peripheral but characterise the
majority of natural language use. The fourfold correspondence is exact for
formal and computational languages and approximate for natural language.

\section{Conclusion}

This paper has traced a fourfold structural correspondence linking formal
grammars, abstract automata, type systems, and proof calculi. The principal
contributions are: (1)~systematic extension of the grammar-automaton
correspondence with stratified precision claims; (2)~identification of the MCS
gap and a concrete conjecture (bounded linear polymorphic calculus);
(3)~cross-linguistic grounding in five typologically diverse languages;
(4)~explicit treatment of gradient grammaticality, multiword expressions, and
the contested status of key empirical arguments; (5)~analysis of the
Merge-application connection; (6)~survey of the neural grammar question; and
(7)~practical implications for programming language design, NLP, and governance
systems.

The convergence of Chomsky's grammars, Turing's machines, Church's types, and
Lambek's categories is one of the deepest structural results of the twentieth
century, and its full articulation remains one of the most compelling open
problems of the twenty-first.

\section*{Acknowledgments}

This work was conducted as part of the Studium Generale ORGANVM research
program. The author acknowledges the use of AI language models (Claude,
Anthropic) as research assistants for literature synthesis, structural
drafting, and adversarial review via the Triadic Review Protocol. All
intellectual claims, arguments, and editorial decisions are the author's own.

\medskip
\noindent\textcopyright~2026 A.~Padavano. This work is licensed under a
Creative Commons Attribution 4.0 International License (CC BY 4.0).

\begin{thebibliography}{99}

\bibitem{Ajdukiewicz1935}
K.~Ajdukiewicz.
\newblock Die syntaktische Konnexitat.
\newblock \emph{Studia Philosophica}, 1:1--27, 1935.

\bibitem{Atkey2018}
R.~Atkey.
\newblock Syntax and semantics of quantitative type theory.
\newblock In \emph{Proc.\ 33rd Annual ACM/IEEE Symp.\ Logic in Computer Science}, pp.~56--65, 2018.

\bibitem{Backus1959}
J.~Backus.
\newblock The syntax and semantics of the proposed international algebraic language of the Zurich ACM-GAMM conference.
\newblock In \emph{Proc.\ Int.\ Conf.\ Information Processing}, UNESCO, pp.~125--131, 1959.

\bibitem{Baker1988}
M.~C.~Baker.
\newblock \emph{Incorporation: A Theory of Grammatical Function Changing}.
\newblock Univ.\ Chicago Press, 1988.

\bibitem{Baker1996}
M.~C.~Baker.
\newblock \emph{The Polysynthesis Parameter}.
\newblock Oxford Univ.\ Press, 1996.

\bibitem{BaldridgeKruijff2003}
J.~Baldridge and G.-J.~M.~Kruijff.
\newblock Multi-modal combinatory categorial grammar.
\newblock In \emph{Proc.\ 10th EACL}, pp.~211--218, 2003.

\bibitem{BarHillel1953}
Y.~Bar-Hillel.
\newblock A quasi-arithmetical notation for syntactic description.
\newblock \emph{Language}, 29(1):47--58, 1953.

\bibitem{Bybee2006}
J.~Bybee.
\newblock From usage to grammar: The mind's response to repetition.
\newblock \emph{Language}, 82(4):711--733, 2006.

\bibitem{Bybee2010}
J.~Bybee.
\newblock \emph{Language, Usage and Cognition}.
\newblock Cambridge Univ.\ Press, 2010.

\bibitem{Chomsky1956}
N.~Chomsky.
\newblock Three models for the description of language.
\newblock \emph{IRE Trans.\ Information Theory}, 2(3):113--124, 1956.

\bibitem{Chomsky1957}
N.~Chomsky.
\newblock \emph{Syntactic Structures}.
\newblock Mouton, 1957.

\bibitem{Chomsky1959}
N.~Chomsky.
\newblock On certain formal properties of grammars.
\newblock \emph{Information and Control}, 2(2):137--167, 1959.

\bibitem{Chomsky1962}
N.~Chomsky.
\newblock Context-free grammars and pushdown storage.
\newblock \emph{MIT Res.\ Lab.\ Electronics Quarterly Progress Report}, 65:187--194, 1962.

\bibitem{Chomsky1995}
N.~Chomsky.
\newblock \emph{The Minimalist Program}.
\newblock MIT Press, 1995.

\bibitem{Chomsky2000}
N.~Chomsky.
\newblock Minimalist inquiries: The framework.
\newblock In R.~Martin et al., eds., \emph{Step by Step}, pp.~89--155. MIT Press, 2000.

\bibitem{Clark2019}
K.~Clark, U.~Khandelwal, O.~Levy, and C.~D.~Manning.
\newblock What does BERT look at? An analysis of BERT's attention.
\newblock In \emph{Proc.\ 2019 ACL Workshop BlackboxNLP}, pp.~276--286, 2019.

\bibitem{ClarkCurran2007}
S.~Clark and J.~R.~Curran.
\newblock Wide-coverage efficient statistical parsing with CCG and log-linear models.
\newblock \emph{Computational Linguistics}, 33(4):493--552, 2007.

\bibitem{CoeckeSadrzadehClark2010}
B.~Coecke, M.~Sadrzadeh, and S.~Clark.
\newblock Mathematical foundations for a compositional distributional model of meaning.
\newblock \emph{Linguistic Analysis}, 36(1--4):345--384, 2010.

\bibitem{Collins2002}
C.~Collins.
\newblock Eliminating labels.
\newblock In S.~D.~Epstein and T.~D.~Seely, eds., \emph{Derivation and Explanation in the Minimalist Program}, pp.~42--64. Blackwell, 2002.

\bibitem{CoquandHuet1988}
T.~Coquand and G.~Huet.
\newblock The Calculus of Constructions.
\newblock \emph{Information and Computation}, 76(2--3):95--120, 1988.

\bibitem{CurryFeys1958}
H.~B.~Curry and R.~Feys.
\newblock \emph{Combinatory Logic}, Vol.~I.
\newblock North-Holland, 1958.

\bibitem{DunfieldKrishnaswami2021}
J.~Dunfield and N.~R.~Krishnaswami.
\newblock Bidirectional typing.
\newblock \emph{ACM Computing Surveys}, 54(5):1--38, 2021.

\bibitem{ErmanWarren2000}
B.~Erman and B.~Warren.
\newblock The idiom principle and the open choice principle.
\newblock \emph{Text}, 20(1):29--62, 2000.

\bibitem{Evey1963}
R.~J.~Evey.
\newblock Application of pushdown store machines.
\newblock \emph{Proc.\ Fall Joint Computer Conference}, pp.~215--227, 1963.

\bibitem{Gazdar1988}
G.~Gazdar.
\newblock Applicability of indexed grammars to natural languages.
\newblock In U.~Reyle and C.~Rohrer, eds., \emph{Natural Language Parsing and Linguistic Theories}, pp.~69--94. Reidel, 1988.

\bibitem{Goldberg1995}
A.~E.~Goldberg.
\newblock \emph{Constructions: A Construction Grammar Approach to Argument Structure}.
\newblock Univ.\ Chicago Press, 1995.

\bibitem{Goldberg2006}
A.~E.~Goldberg.
\newblock \emph{Constructions at Work}.
\newblock Oxford Univ.\ Press, 2006.

\bibitem{Hale1983}
K.~Hale.
\newblock Warlpiri and the grammar of non-configurational languages.
\newblock \emph{Natural Language and Linguistic Theory}, 1(1):5--47, 1983.

\bibitem{HewittLiang2019}
J.~Hewitt and P.~Liang.
\newblock Designing and interpreting probes with control tasks.
\newblock In \emph{Proc.\ EMNLP 2019}, pp.~2733--2743, 2019.

\bibitem{HewittManning2019}
J.~Hewitt and C.~D.~Manning.
\newblock A structural probe for finding syntax in word representations.
\newblock In \emph{Proc.\ NAACL-HLT 2019}, pp.~4129--4138, 2019.

\bibitem{Howard1980}
W.~A.~Howard.
\newblock The formulae-as-types notion of construction.
\newblock In J.~P.~Seldin and J.~R.~Hindley, eds., \emph{To H.~B.~Curry}, pp.~479--490. Academic Press, 1980. (Manuscript 1969.)

\bibitem{Jawahar2019}
G.~Jawahar, B.~Sagot, and D.~Seddah.
\newblock What does BERT learn about the structure of language?
\newblock In \emph{Proc.\ 57th ACL}, pp.~3651--3657, 2019.

\bibitem{Joshi1985}
A.~K.~Joshi.
\newblock Tree adjoining grammars: How much context-sensitivity is required to provide reasonable structural descriptions?
\newblock In D.~R.~Dowty et al., eds., \emph{Natural Language Parsing}, pp.~206--250. Cambridge Univ.\ Press, 1985.

\bibitem{JoshiLevyTakahashi1975}
A.~K.~Joshi, L.~S.~Levy, and M.~Takahashi.
\newblock Tree adjunct grammars.
\newblock \emph{J.\ Computer and System Sciences}, 10(1):136--163, 1975.

\bibitem{KaplanKay1994}
R.~Kaplan and M.~Kay.
\newblock Regular models of phonological rule systems.
\newblock \emph{Computational Linguistics}, 20(3):331--378, 1994.

\bibitem{Kartsaklis2021}
D.~Kartsaklis, I.~Fan, R.~Yeung, A.~Pearson, R.~Lorenz, A.~Toumi, G.~de~Felice, K.~Meichanetzidis, S.~Clark, and B.~Coecke.
\newblock lambeq: An efficient high-level Python library for quantum NLP.
\newblock arXiv:2110.04236, 2021.

\bibitem{Kleene1956}
S.~C.~Kleene.
\newblock Representation of events in nerve nets and finite automata.
\newblock In C.~E.~Shannon and J.~McCarthy, eds., \emph{Automata Studies}, pp.~3--42. Princeton Univ.\ Press, 1956.

\bibitem{Kuroda1964}
S.-Y.~Kuroda.
\newblock Classes of languages and linear-bounded automata.
\newblock \emph{Information and Control}, 7(2):207--223, 1964.

\bibitem{Lambek1958}
J.~Lambek.
\newblock The mathematics of sentence structure.
\newblock \emph{American Mathematical Monthly}, 65(3):154--170, 1958.

\bibitem{Lambek1972}
J.~Lambek.
\newblock Deductive systems and categories III.
\newblock In F.~W.~Lawvere, ed., \emph{Toposes, Algebraic Geometry and Logic}, LNCS~274, pp.~57--82. Springer, 1972.

\bibitem{Lambek1980}
J.~Lambek.
\newblock From lambda calculus to Cartesian closed categories.
\newblock In J.~P.~Seldin and J.~R.~Hindley, eds., \emph{To H.~B.~Curry}, pp.~375--402. Academic Press, 1980.

\bibitem{Lambek1999}
J.~Lambek.
\newblock Type grammar revisited.
\newblock In A.~Lecomte et al., eds., \emph{Logical Aspects of Computational Linguistics}, LNCS~1582, pp.~1--27. Springer, 1999.

\bibitem{Lambek2008}
J.~Lambek.
\newblock \emph{From Word to Sentence}.
\newblock Polimetrica, 2008.

\bibitem{Langacker1987}
R.~W.~Langacker.
\newblock \emph{Foundations of Cognitive Grammar}, Vol.~I.
\newblock Stanford Univ.\ Press, 1987.

\bibitem{Langacker2008}
R.~W.~Langacker.
\newblock \emph{Cognitive Grammar: A Basic Introduction}.
\newblock Oxford Univ.\ Press, 2008.

\bibitem{LauClarkLappin2017}
J.~H.~Lau, A.~Clark, and S.~Lappin.
\newblock Grammaticality, acceptability, and probability.
\newblock \emph{Cognitive Science}, 41(5):1202--1241, 2017.

\bibitem{LewisSteedman2014}
M.~Lewis and M.~Steedman.
\newblock A* CCG parsing with a supertag-factored model.
\newblock In \emph{Proc.\ EMNLP 2014}, pp.~990--1000, 2014.

\bibitem{MartinLof1975}
P.~Martin-L\"{o}f.
\newblock An intuitionistic theory of types: Predicative part.
\newblock In H.~E.~Rose and J.~C.~Shepherdson, eds., \emph{Logic Colloquium '73}, pp.~73--118. North-Holland, 1975.

\bibitem{MartinLof1984}
P.~Martin-L\"{o}f.
\newblock \emph{Intuitionistic Type Theory}.
\newblock Bibliopolis, 1984.

\bibitem{McBride2016}
C.~McBride.
\newblock I got plenty o' nuttin'.
\newblock In S.~Lindley et al., eds., \emph{A List of Successes That Can Change the World}, LNCS~9600, pp.~207--233. Springer, 2016.

\bibitem{MichaelisKracht1997}
J.~Michaelis and M.~Kracht.
\newblock Semilinearity as a syntactic invariant.
\newblock In C.~Retor\'{e}, ed., \emph{Logical Aspects of Computational Linguistics}, LNCS~1328, pp.~329--345. Springer, 1997.

\bibitem{Montague1970a}
R.~Montague.
\newblock Universal grammar.
\newblock \emph{Theoria}, 36(3):373--398, 1970.

\bibitem{Montague1970b}
R.~Montague.
\newblock English as a formal language.
\newblock In B.~Visentini et al., eds., \emph{Linguaggi nella Societa e nella Tecnica}, pp.~189--224. Edizioni di Comunita, 1970.

\bibitem{Montague1973}
R.~Montague.
\newblock The proper treatment of quantification in ordinary English.
\newblock In K.~J.~J.~Hintikka et al., eds., \emph{Approaches to Natural Language}, pp.~221--242. Reidel, 1973.

\bibitem{Myhill1957}
J.~Myhill.
\newblock Finite automata and the representation of events.
\newblock WADD Technical Report 57--624, 1957.

\bibitem{Naur1960}
P.~Naur, ed.
\newblock Report on the algorithmic language ALGOL 60.
\newblock \emph{Communications of the ACM}, 3(5):299--314, 1960.

\bibitem{Pentus1997}
M.~Pentus.
\newblock Product-free Lambek calculus and context-free grammars.
\newblock \emph{J.\ Symbolic Logic}, 62(2):648--660, 1997.

\bibitem{Pollard1984}
C.~Pollard.
\newblock \emph{Generalized Phrase Structure Grammars, Head Grammars, and Natural Language}.
\newblock PhD thesis, Stanford University, 1984.

\bibitem{PullumGazdar1982}
G.~K.~Pullum and G.~Gazdar.
\newblock Natural languages and context free languages.
\newblock \emph{Linguistics and Philosophy}, 4(4):471--504, 1982.

\bibitem{RabinScott1959}
M.~O.~Rabin and D.~Scott.
\newblock Finite automata and their decision problems.
\newblock \emph{IBM J.\ Research and Development}, 3(2):114--125, 1959.

\bibitem{SandlerLilloMartin2006}
W.~Sandler and D.~Lillo-Martin.
\newblock \emph{Sign Language and Linguistic Universals}.
\newblock Cambridge Univ.\ Press, 2006.

\bibitem{Shieber1985}
S.~M.~Shieber.
\newblock Evidence against the context-freeness of natural language.
\newblock \emph{Linguistics and Philosophy}, 8(3):333--343, 1985.

\bibitem{SprouseSchutzeAlmeida2013}
J.~Sprouse, C.~T.~Sch\"{u}tze, and D.~Almeida.
\newblock A comparison of informal and formal acceptability judgments.
\newblock \emph{Lingua}, 134:34--44, 2013.

\bibitem{Stabler1997}
E.~P.~Stabler.
\newblock Derivational minimalism.
\newblock In C.~Retor\'{e}, ed., \emph{Logical Aspects of Computational Linguistics}, LNCS~1328, pp.~68--95. Springer, 1997.

\bibitem{Stabler2011}
E.~P.~Stabler.
\newblock Computational perspectives on minimalism.
\newblock In C.~Boeckx, ed., \emph{The Oxford Handbook of Linguistic Minimalism}, pp.~617--641. Oxford Univ.\ Press, 2011.

\bibitem{Stabler2013}
E.~P.~Stabler.
\newblock Two models of minimalist, incremental syntactic analysis.
\newblock \emph{Topics in Cognitive Science}, 5(3):611--633, 2013.

\bibitem{Steedman2000}
M.~Steedman.
\newblock \emph{The Syntactic Process}.
\newblock MIT Press, 2000.

\bibitem{VijayShankerWeir1994}
K.~Vijay-Shanker and D.~J.~Weir.
\newblock The equivalence of four extensions of context-free grammars.
\newblock \emph{Mathematical Systems Theory}, 27(6):511--546, 1994.

\bibitem{Wray2002}
A.~Wray.
\newblock \emph{Formulaic Language and the Lexicon}.
\newblock Cambridge Univ.\ Press, 2002.

\end{thebibliography}

\end{document}
